\section{特征函数怎样证明中心极限定理}
\label{sec:05-Probability/08-Limit-Theorems/04}

上一节留下一个刺眼的事实：起点可以是均匀分布、Bernoulli 分布或指数分布，标准化之后的极限却是同一条 \(e^{-x^2/2}/\sqrt{2\pi}\)。原始分布的一切细节——离散还是连续、对称还是偏斜、支撑在哪里——全部消失了。

大数定律不需要解释这种现象，它只要把方差摊薄就完事。中心极限定理不一样，它断言的是整个分布形状的极限，而形状不是某一个矩，压不下去。

先看看正面强攻有多难。取 \(X_1,\ldots,X_n\) 独立且服从 \(\Uniform(0,1)\)，问 \(S_n=X_1+\cdots+X_n\) 的密度。两个的和是三角形，三个的和是分段二次，\(n\) 个的和是有 \(n\) 段的 \(n-1\) 次多项式，分段点越来越密。换一个母分布，连这个分段多项式都写不出来。独立和的密度是卷积（见\hyperref[sec:05-Probability/05-Joint-Distributions/03]{``随机变量之和与卷积''}），而卷积序列的极限，直接算是算不动的。

所以要绕路。

\begin{center}
\begin{tikzpicture}[x=1cm,y=1cm,font=\small,>=stealth]
  \draw[black!45,fill=blue!8] (-2.7,2.15) rectangle (2.7,3.05);
  \node[align=center] at (0,2.6) {标准化和 \(Z_n\) 的分布\\（\(n\) 重卷积）};
  \draw[black!45,fill=green!10] (5.9,2.15) rectangle (11.5,3.05);
  \node at (8.7,2.6) {\(\varphi_{Z_n}(t)=\bigl[\varphi_Y(t/\sqrt n)\bigr]^n\)};
  \draw[black!45,fill=green!10] (7.4,-0.45) rectangle (10.0,0.45);
  \node at (8.7,0) {\(e^{-t^2/2}\)};
  \draw[black!45,fill=orange!14] (-1.6,-0.45) rectangle (1.6,0.45);
  \node at (0,0) {\(\mathcal N(0,1)\)};
  \draw[->,very thick] (2.75,2.6) -- (5.85,2.6);
  \node[above,black!70,font=\footnotesize] at (4.3,2.68) {取特征函数：卷积变乘法};
  \draw[->,very thick] (8.7,2.1) -- (8.7,0.5);
  \node[right,black!70,font=\footnotesize,align=left] at (8.85,1.3) {原点处展开到二阶，\\再用 \((1+x/n)^n\to e^x\)};
  \draw[->,very thick] (7.35,0) -- (1.65,0);
  \node[above,black!70,font=\footnotesize] at (4.5,0.08) {Lévy 连续性定理};
  \draw[dashed,black!55] (0,2.1) -- (0,0.5);
  \draw[black!70,thick] (-0.22,1.12) -- (0.22,1.48);
  \draw[black!70,thick] (-0.22,1.48) -- (0.22,1.12);
  \node[left,black!70,font=\footnotesize,align=right] at (-0.35,1.3) {对卷积序列\\直接取极限：\\走不通};
\end{tikzpicture}

\emph{左边那条竖线是被划掉的直路。整个证明就是这张图上那条绕行的路：换一种表示让求和变简单，在新表示里算极限，再把结果翻译回分布。走完之后你会发现，难点从``分析卷积序列''降到了``算一个 \((1-c/n)^n\) 的极限''。}
\end{center}

\subsection{为什么是特征函数，而不是矩母函数}

特征函数把分布编码成
\[
\varphi_X(t)=\E[e^{itX}],
\]
关键性质是独立变量相加对应编码相乘：\(\varphi_{X+Y}(t)=\E[e^{itX}]\E[e^{itY}]=\varphi_X(t)\varphi_Y(t)\)。卷积就此消失。

矩母函数 \(M_X(s)=\E[e^{sX}]\) 也有同样的乘法性质，用起来还不必碰复数，为什么不用它？因为它可能不存在。被积函数 \(e^{sx}\) 在 \(x\to\infty\) 时爆炸式增长，只要分布的右尾比指数厚一点，积分就发散：对数正态分布的矩母函数在任何 \(s>0\) 处都是无穷，Cauchy 分布在任何 \(s\neq0\) 处都是无穷。而特征函数的被积函数模长恒为 \(1\)，积分对任何分布都绝对收敛。

这个差别不是洁癖。要证明一条对所有有限方差分布都成立的定理，工具本身就得对所有分布都有定义，否则每次用之前还要先检查前提，而且恰恰是最需要警惕的那些重尾分布会把工具卡住。矩母函数在集中不等式里很好用（见\hyperref[sec:05-Probability/07-Transforms-and-Bounds/04]{``Chernoff 方法与 Hoeffding 集中界''}），因为那里本来就假设了轻尾；在极限定理里它不够格。

还需要两条已经建立过的事实（见\hyperref[sec:05-Probability/07-Transforms-and-Bounds/01]{``生成函数、矩母函数与特征函数''}）：特征函数唯一确定分布；它在原点附近的各阶导数给出各阶矩，\(\varphi^{(k)}(0)=i^k\E[X^k]\)。前者保证算出极限之后能认人，后者保证展开时能把矩塞进去。

最后一块拼图是 Lévy 连续性定理：\(X_n\) 依分布收敛到 \(X\)，当且仅当 \(\varphi_{X_n}(t)\to\varphi_X(t)\) 对每个实数 \(t\) 成立。有了它，``分布形状接近''这件难以直接操作的事，被翻译成了逐点收敛这件普通的分析工作。

\subsection{四步骨架}

设 \(X_i\) 独立同分布，\(\E X_i=\mu\)，\(\Var(X_i)=\sigma^2<\infty\)。令 \(Y_i=(X_i-\mu)/\sigma\)，于是 \(\E Y_i=0\)、\(\E[Y_i^2]=1\)，标准化和为 \(Z_n=n^{-1/2}\sum_i Y_i\)。

第一步，乘积出现。独立性让特征函数相乘，同分布让所有因子相同：
\[
\varphi_{Z_n}(t)=\Bigl[\varphi_Y\bigl(t/\sqrt n\bigr)\Bigr]^n .
\]

第二步，在原点展开到二阶。用 \(\varphi_Y(0)=1\)、\(\varphi_Y'(0)=i\E Y=0\)、\(\varphi_Y''(0)=-\E[Y^2]=-1\)，
\[
\varphi_Y(s)=1-\frac{s^2}2+o(s^2),\qquad s\to0 .
\]
一阶项因为中心化而消失，二阶项由方差归一而定为 \(-1/2\)。代入 \(s=t/\sqrt n\)，得
\[
\varphi_Y\bigl(t/\sqrt n\bigr)=1-\frac{t^2}{2n}+o\!\left(\frac1n\right).
\]

第三步，取 \(n\) 次幂。这就是微积分里的 \((1+x_n/n)^n\to e^x\)：
\[
\varphi_{Z_n}(t)=\left[1-\frac{t^2/2+o(1)}{n}\right]^n\longrightarrow e^{-t^2/2}.
\]

第四步，认人。\(e^{-t^2/2}\) 正是标准正态的特征函数，Lévy 连续性定理于是给出 \(Z_n\xrightarrow{d}\mathcal N(0,1)\)。

余项的逐阶控制、\(o(\cdot)\) 在取 \(n\) 次幂时的严格处理、复对数取主值的分支问题，这些都要写清楚才算完整证明，但它们不改变主意。主意只有一句：卷积换成乘法之后，剩下的是大一微积分。

值得停下来看一眼极限里剩了什么。原始分布的信息一共只用掉了两个数——均值和方差，而它们在标准化那一步就被规范成了 \(0\) 和 \(1\)。三阶以上的矩全被 \(o(s^2)\) 吞掉了。Gaussian 的普适性来源于此：不是所有分布最后都长得像正态，而是标准化只留下前两阶信息，前两阶相同的东西极限自然相同。

\subsection{每一步的前提各挡住什么}

独立性挡住第一步的失败。如果只假设两两不相关，和的分布不由边缘分布决定，特征函数写不成乘积，整条路从起点就断了。

有限方差挡住第二步的失败。Cauchy 分布的特征函数是 \(e^{-\abs t}\)，它在 \(t=0\) 处连一阶导数都没有，谈不上展开到二阶。这不是技术障碍，是实情：同样的``相乘再缩放''操作作用在 \(e^{-\abs t}\) 上得到 \((e^{-\abs t/n})^n=e^{-\abs t}\)，特征函数原地不动，样本平均永远是 Cauchy。有限方差画出的是疆界，疆界外面是稳定分布的地盘。

同分布挡住的只是简洁性，不是正确性。分布各异时乘积不再是同一因子的 \(n\) 次幂，需要逐项控制每个因子的展开，保证没有哪一项独自主导。上一节的 Lindeberg 条件正是这件事的精确版本。

顺带说，Berry--Esseen 的入口也在这里：把展开做到三阶并认真估计余项随 \(n\) 衰减的速度，就得到有限 \(n\) 的误差界，而不只是极限陈述。同一个 Taylor 展开，区别在于这次不把高阶项扔掉，而是把它们捆起来算账。

\subsection{同一台机器的其他产出}

写成幂、展开、取极限、认人——这套流程不止能证一条定理。Binomial\((n,\lambda/n)\) 的特征函数是 \((1+\lambda(e^{it}-1)/n)^n\)，取极限得到 \(\exp(\lambda(e^{it}-1))\)，正是 Poisson\((\lambda)\) 的特征函数，于是``大量试验中的稀有事件趋向 Poisson''得到了严格证明。若只展开到一阶，\(\varphi_X(u)=1+i\mu u+o(u)\)，代入 \(u=t/n\) 得 \(\varphi_{\bar X_n}(t)\to e^{it\mu}\)，这是常数 \(\mu\) 的特征函数——弱大数定律。方差不存在而尾指数为 \(\alpha\in(0,2)\) 时，原点附近的展开带有 \(\abs t^\alpha\) 项，极限变成 \(e^{-c\abs t^\alpha}\)，这就是稳定分布族。

三条定理的差别只在于展开到哪一阶。大数定律用一阶，中心极限定理用二阶，稳定律用分数阶。以后遇到任何``某分布收敛到某分布''的命题，第一个问题都可以是：它的特征函数在原点附近能展开到几阶。

\subsection{这台机器不保证的事}

连续性定理要求对每个固定的 \(t\) 逐点收敛，检查一个 \(t\) 不算数；反过来它也不要求一致收敛或导数收敛，逐点就够。

更微妙的是，逐点极限必须是某个合法分布的特征函数，而这一条不能想当然。取 \(X_n\sim\Uniform(-n,n)\)，特征函数 \(\varphi_{X_n}(t)=\sin(nt)/(nt)\) 逐点趋于一个在 \(t=0\) 取 \(1\)、其余处取 \(0\) 的函数。特征函数必须处处连续，这个极限不是任何分布的特征函数——事实上概率质量跑到无穷远去了，分布根本不收敛。检查极限函数在 \(t=0\) 处是否连续，就是在检查质量有没有逃逸。

还有一点已经强调过但值得重复：特征函数证明给的是极限陈述，不给有限 \(n\) 的误差界。上一节那张右尾对数图上的七倍差距，不会因为你把这个证明读懂而缩小，它需要 Berry--Esseen 或 Edgeworth 展开单独估计。证明告诉你极限存在且是谁，不告诉你现在离它多远。

到这里，本单元的两条主结论都有了完整交代：均值会稳定，波动按 \(1/\sqrt n\) 收缩且形状渐近正态。下一节把这两条反过来用——不再是有了模型去预测数据的行为，而是造出随机数据来算一个确定的量。误差按 \(1/\sqrt n\) 衰减这件事，在那里会从一个理论性质变成一条报价单。
